\documentclass[12pt]{article}
\usepackage[english]{babel}
\usepackage{graphicx}
\usepackage{framed}
\usepackage[normalem]{ulem}
\usepackage{amsmath}
\usepackage{amsthm}
\usepackage{amssymb}
\usepackage{amsfonts}
\usepackage{enumerate}
\usepackage{mathtools}
\usepackage[utf8]{inputenc}
\usepackage[margin=1in]{geometry}
\usepackage{float} % for "H" placement specifier
\usepackage[authoryear]{natbib}
\usepackage{indentfirst}
\usepackage[dvipsnames]{xcolor}
\usepackage{array}
\usepackage{subcaption}
\usepackage{caption}

\makeatletter
\newcommand*{\addFileDependency}[1]{% argument=file name and extension
  \typeout{(#1)}
  \@addtofilelist{#1}
  \IfFileExists{#1}{}{\typeout{No file #1.}}
}
\makeatother

\newcommand*{\myexternaldocument}[1]{%
    \externaldocument{#1}%
    \addFileDependency{#1.tex}%
    \addFileDependency{#1.aux}%
}

% ==== Cross-document referencing (MUST be before hyperref) ====
\usepackage{xr-hyper}

% Tell LaTeX which external document to read labels from
% Make sure eck2023draft.aux is in the SAME folder
\myexternaldocument{yan2026paradoxes}
% ==== Hyperref comes AFTER xr-hyper ====
\usepackage{hyperref}
\hypersetup{
  linkcolor  = blue,
  citecolor  = blue,
  urlcolor   = blue,
  colorlinks = true,
}

\usepackage[T1]{fontenc}
\usepackage{tikz}
\newcommand{\fatterdot}{\raisebox{0.15ex}{\tikz\filldraw[black,x=1.5pt,y=1.5pt] (0,0) circle (1);}}
\newcommand{\xn}{x^{\frac{1}{n}}}
\newcommand{\xnone}{x^{\frac{1}{n+1}}}

\newcommand{\R}{\mathbb{R}}
\DeclareUnicodeCharacter{FB01}{fi}
\DeclareUnicodeCharacter{FF0C}{,}
\DeclareMathOperator*{\Prob}{P}
\DeclareMathOperator*{\argmax}{arg\,max}
\DeclareMathOperator*{\argmin}{arg\,min}
\setlength\extrarowheight{1pt}

\newcommand\red[1]{{\color{red}#1}}

\bibliographystyle{plainnat}

\newtheorem{theorem}{Theorem}[section]
\newtheorem{corollary}{Corollary}[theorem]
\newtheorem{lemma}[theorem]{Lemma}
\newtheorem{prop}[theorem]{Proposition}
\newtheorem{defn}[theorem]{Definition}
\allowdisplaybreaks

\usepackage{setspace}
\doublespacing

% Option B (comment A above and uncomment B below): biblatex
%\usepackage[backend=biber,style=authoryear,sorting=nyt]{biblatex}
%\addbibresource{supp.bib}

% ---------- Title block ----------
\title{Supplementary Materials for ``A Formal Investigation of Paradoxes Between Latent Aptitude and Observed Achievement via Full House Methodology''}
\date{}

% Optional footnote-style affiliations
\newcommand{\affil}[2]{\par\noindent\textsuperscript{#1}\,\textit{#2}}

\begin{document}
\maketitle



This supplementary materials includes proofs of technical results and additional numerical examples supporting conclusions from Section 4 the main manuscript. 

\tableofcontents

%\section{Monotonicity}
%
%\begin{lemma}\label{lemma:monotonicity}
%Let $1 \leq j < n$ and let $x \in (0,1)$. Let $I_{\alpha, \beta}(x)$ be a regularized incomplete beta function as in Definition~\ref{def:ribf}. Then, $I_{n-j, j+1}(\xn) \geq I_{n-j+1, j+1}(\xnone)$.
%\end{lemma}
%%
%%\begin{center}
%%	\textbf{Need to fix notation; should be $I_{\alpha,\beta}(x)$}
%%\end{center}
%
%\begin{proof}
%We can define a function $f_{n,j}(x)$ as
%\begin{align*}
%  f_{n,j}(x) & = I_{n-j, j+1}(\xn) - I_{n-j+1, j+1}(\xnone) \\
%  &= \int_0^{\xn}\frac{y^{n-j-1}(1-y)^j}{B(n-j,j+1)}dy 
%    - \int_0^{\xnone}\frac{y^{n-j}(1-y)^j}{B(n-j+1,j+1)}dy.
%\end{align*}
%An application of the chain rule yields 
%\begin{align*}
%  &f_{n,j}^\prime(x) = \frac{\left(\xn\right)^{n-j-1}\left(1-\xn\right)^j\left(x^{\frac{1-n}{n}}\right)}{nB(n-j,j+1)} 
%    - \frac{\left(\xnone\right)^{n-j}\left(1-\xnone\right)^j\left(x^{-\frac{n}{n+1}}\right)}{(n+1)B(n-j+1,j+1)} \\
%  &= \frac{\left(\xn\right)^{n-j-1}\left(1-\xn\right)^j\left(x^{\frac{1-n}{n}}\right)}{nB(n-j,j+1)} 
%    - \frac{\left(\xnone\right)^{n-j}\left(1-\xnone\right)^j\left(x^{-\frac{n}{n+1}}\right)}{(n-j)B(n-j,j+1)} \\    
%  &= \frac{1}{B(n-j,j+1)}\left(\frac{\left(\xn\right)^{n-j-1}\left(1-\xn\right)^j\left(x^{\frac{1-n}{n}}\right)}{n} 
%    - \frac{\left(\xnone\right)^{n-j}\left(1-\xnone\right)^j\left(x^{-\frac{n}{n+1}}\right)}{n-j}\right),
%\end{align*}
%where we used $B(n-j+1,j+1) = \frac{n-j}{n+1}B(n-j,j+1)$ in the second equality. Now define $A_{n,j} = \{x:f_{n,j}^\prime(x) \geq 0\}$. For $x \in A_{n,j}$ we have
%\begin{align*}
%	&\frac{\left(\xn\right)^{n-j-1}\left(1-\xn\right)^j\left(x^{\frac{1-n}{n}}\right)}{n} 
%    \geq \frac{\left(\xnone\right)^{n-j}\left(1-\xnone\right)^j\left(x^{-\frac{n}{n+1}}\right)}{n-j} \\
%    &\qquad\implies \frac{\left(\xn\right)^{n-j-1}\left(1-\xn\right)^j\left(x^{\frac{1-n}{n}}\right)}{\left(\xnone\right)^{n-j}\left(1-\xnone\right)^j\left(x^{-\frac{n}{n+1}}\right)} \geq \frac{n}{n-j} \\
%    &\qquad\implies x^{-\frac{j}{n(n+1)}}\left(\frac{1-\xn}{1-\xnone}\right)^j \geq \frac{n}{n-j} \\
%    &\qquad\implies \frac{x^{-\frac{1}{n(n+1)}}-\xnone}{1-\xnone} \geq \left(\frac{n}{n-j}\right)^\frac{1}{j}.
%\end{align*}
%Define $g_n(x) = \frac{x^{-\frac{1}{n(n+1)}}-\xnone}{1-\xnone}$. Then 
%\begin{align*}
%	g_n(x) = \frac{x^{-\frac{1}{n(n+1)}}-\xnone}{1-\xnone} = \frac{x^{-\frac{1}{n(n+1)}} - 1}{1-\xnone} + 1,
%\end{align*}
%and
%\begin{align*}
%	g^{'}_n(x) &= \frac{1}{(1-\xnone)^2}\left(-\frac{1}{n(n+1)}x^{-\frac{1}{n(n+1)}-1}(1-\xnone) - (x^{-\frac{1}{n(n+1)}} - 1)(-\frac{1}{n+1}x^{\frac{1}{n+1}-1})\right)\\
%	&= \frac{1}{(1-\xnone)^2}\left( \frac{1}{n}x^{-\frac{n^2+1}{n(n+1)}} - \frac{1}{n+1}x^{-\frac{n}{n+1}} - \frac{1}{n(n+1)}x^{-\frac{n^2 + n + 1}{n(n+1)}}\right)\\
%	&= \frac{x^{-\frac{n^2+1}{n(n+1)}}}{n(n+1)(1-\xnone)^2}\left((n+1) - nx^{\frac{1}{n(n+1)}} - x^{-\frac{1}{n+1}}\right).
%\end{align*}
%Now define $l_n(x) = (n+1) - nx^{\frac{1}{n(n+1)}} - x^{-\frac{1}{n+1}}$ and see that
%\begin{align*}
%	l^{'}_n(x) = -\frac{1}{n+1}x^{\frac{1}{n(n+1)} - 1} + \frac{1}{n+1}x^{-\frac{1}{n+1} - 1} = \frac{1}{n+1}x^{-\frac{n+2}{n+1}}\left( 1 - \xn \right).
%\end{align*}
%Thus, $l_n(x)$ is an increasing function for $x \in (0,1)$ and $\sup_{x \in (0,1)}l_n(x) = 0$. Therefore, $l_n(x) < 0$ for $x \in (0,1)$. Thus, $g_n(x)$ is an decreasing function, and $g_n(x)\to1$ as $x\to1$. We can see that $g_n(x) \in (1,+\infty)$ for $x \in (0,1)$. Thus $A_{n,j} = (0,x^*_{n,j}]$ where $x^*_{n,j}$ is such that $g_n(x^*_{n,j}) = (\frac{n}{n-j})^j$. Putting it all together, $f_{n,j}(0) = f_{n,j}(1) = 0$ and this implies that $I_{n-j, j+1}(\xn) \geq I_{n-j+1, j+1}(\xnone)$ for all $x \in (0,1)$.
%\end{proof}


\section{Proofs}


\subsection{Proof of Theroem 2.2}
\begin{proof}
Fix a population $i$ and suppress the index $i$. First, since $F_X$ is continuous, we can express well-defined order statistics as $X_{(1)} < \cdots < X_{(N)}$. By rank-based selection, the observed system consists of indices $j$ such that $r_j \in \mathcal{R}$, yielding observed pairs $(Y_{(j)}, X_{(r_j)})$, for $j = 1,\dots,n$. By assumption, for each observed rank $j$ there exists a strictly increasing measurable function $\psi_j$ such that
$$
  Y_{(j)} = \psi_j(X_{(r_j)}) \quad \text{a.s.}
$$

We now identify the distribution function $F_{Y_{(j)}}$ for fixed $j$. Since $\psi_j$ is strictly increasing, it admits a measurable inverse on its range. For any $y \in \mathbb{R}$,
$$
\mathbb{P}(Y_{(j)} \le y) 
	= \mathbb{P}(\psi_j(X_{(r_j)}) \le y)
    = \mathbb{P}(X_{(r_j)} \le \psi_j^{-1}(y)).
$$
Hence, $F_{Y_{(j)}}(y) = F_{X_{(r_j)}}(\psi_j^{-1}(y))$, which implies that
$$
  F_{Y_{(j)}}(\psi_j(x)) = F_{X_{(r_j)}}(x)
	\quad \text{for } F_{X_{(r_j)}}\text{-a.e. } x.
$$

We now show that $Y_{(j)}$ is recovered from the Full House Modeling construction. Since $F_X$ is continuous, the distribution of the order statistic $X_{(r_j)}$ is also continuous and strictly increasing on the interior of its support, hence admits a well-defined generalized inverse. Likewise, $F_{Y_{(j)}}$ is continuous and strictly increasing. Thus both admit well-defined inverses. Let $u = F_{X_{(r_j)}}(x)$. Then $F_{Y_{(j)}}(\psi_j(x)) = u$, so that 
$$
\psi_j(x)
	= F_{Y_{(j)}}^{-1}(u)
	= F_{Y_{(j)}}^{-1}\!\big(F_{X_{(r_j)}}(x)\big)
$$
for $F_{X_{(r_j)}}$-almost every $x$. Therefore,
\begin{equation} \label{eq:FHM}
Y_{(j)} = F_{Y_{(j)}}^{-1} \circ F_{X_{(r_j)}}(X_{(r_j)})
	\quad \text{a.s.}	
\end{equation}

We now verify the uniqueness of \eqref{eq:FHM}. Let $U = F_{X_{(r_j)}}(X_{(r_j)}) \sim \mathrm{Unif}(0,1)$ by the probability integral transform, and define
$$
  g(u) := \psi_j\!\big(F_{X_{(r_j)}}^{-1}(u)\big).
$$
Then $g$ is strictly increasing and $g(U) \sim F_{Y_{(j)}}$. By the fact that if $U\sim\mathrm{Unif}(0,1)$ and $g$ is increasing, then the distribution of $g(U)$ uniquely determines $g$ almost everywhere, implying $g(u)=F^{-1}_{Y_{(j)}}(u)$ a.e. It follows that 
$$
\psi_j(x) = F_{Y_{(j)}}^{-1}\!\big(F_{X_{(r_j)}}(x)\big)
	\quad \text{for } F_{X_{(r_j)}}\text{-a.e. } x.
$$
Thus, the strictly increasing coupling is unique up to null sets. The representation holds for each observed rank $j$, and uniqueness implies that the joint distribution of $(X,Y)$ coincides with that induced by the Full House Model.
\end{proof}






%%%% Proof of Lemma 3.2
\subsection{Proof of Lemma~\ref{lemma:rep}}

\begin{proof}
Define $f_{1,j}(\cdot) = F_{Y_{1, (j)}}^{-1} \circ F_{X_{1,(r_{1,j})}} (\cdot)$ and 
$f_{2,k}(\cdot) = F_{Y_{2, (k)}}^{-1} \circ F_{X_{2,(r_{2,k})}} (\cdot)$. Then we have 
\begin{align*}
  &\Prob\left(X_{1,(r_{1,j})} > X_{2,(r_{2,k})}, Y_{1,(j)} < Y_{2,(k)}\right) \\
  &\qquad= \Prob\left(X_{1,(r_{1,j})} > X_{2,(r_{2,k}))}, f_{1,j}(X_{1,(r_{1,j})}) < f_{2,k}(X_{2,(r_{2,k})})\right) \\
  &\qquad= \Prob\left( X_{2,(r_{2,k})} < X_{1,(r_{1,j})}  < h_{j,k}\left(X_{2,(r_{2,k})}\right)\right),
\end{align*}
where $h_{j,k}(\cdot) = (f_{1,j}^{-1} \circ f_{2,k})(\cdot)$. %Existence of $f_i$ follows directly from the existence of $F_Y^{-1}(y)$, and existence of $h_{i,j}$ follows directly from the existence of $f_i$ and the existence of $F_X^{-1}(x)$.
A similar technique can be applied to the second equation to obtain the result.
\end{proof}


%%%% Proof of Theorem 3.4
\subsection{Proof of Theorem~\ref{theorem:samedistgeneral}} 

\begin{proof}
By Lemma~\ref{lemma:rep},
$$
	\Prob\left(X_{1,(r_{1,j})} > X_{2,(r_{2,k})}, Y_{1,(j)} < Y_{2,(k)}\right) = \Prob\left( X_{2,(r_{2,k})} < X_{1,(r_{1,j})}  < h_{j,k}\left(X_{2,(r_{2,k})}\right)\right),
$$
where $h_{j,k}(\cdot) = F_{X_{1,(r_{1,j})}}^{-1} \circ F_{Y_{1, (j)}} \circ F_{Y_{2, (k)}}^{-1} \circ F_{X_{2,(r_{2,k})}} (\cdot)$. Under the assumption $F_{Y_1}=F_{Y_2}=F_Y$, the distributions of the order statistics admit Beta representations, which allow us to express the mapping $h_{j,k}$ in terms of regularized incomplete beta functions.
\begin{align*}
	&\Prob\left( X_{2,(r_{2,k})} < X_{1,(r_{1,j})}  < h_{j,k}\left(X_{2,(r_{2,k})}\right)\right) \\
	&= \int_{-\infty}^{+\infty}\left\{\int_{x}^{F_X^{-1} \circ I_{r_{1,j}, N_1-r_{1,j}+1}^{-1} \circ I_{j, n_1-j+1} \circ I_{k, n_2-k+1}^{-1} \circ I_{r_{2,k}, N_2-r_{2,k}+1} \circ F_X(x)} f_{X_{1,(r_{1,j})}}(y)\ dy\right\} \ f_{X_{2,(r_{2,k})}}(x) \ dx \\
	&= \int_{-\infty}^{+\infty}\left\{\int_{x}^{F_X^{-1} \circ I_{r_{1,j}, N_1-r_{1,j}+1}^{-1} \circ I_{j, n_1-j+1} \circ I_{k, n_2-k+1}^{-1} \circ I_{r_{2,k}, N_2-r_{2,k}+1} \circ F_X(x)} d \ I_{ r_{1,j}, N_1- r_{1,j}+1}(F_X(y))\right\} \ f_{X_{2,(r_{2,k})}}(x) \ dx \\
	&= \int_{-\infty}^{+\infty}\left\{I_{j, n_1-j+1} \circ I_{k, n_2-k+1}^{-1} \circ I_{ r_{2,k}, N_2- r_{2,k}+1} \circ F_X(x) - I_{ r_{1,j}, N_1- r_{1,j}+1} \circ F_X(x)\right\}^+ \ f_{X_{2,(r_{2,k})}}(x) \ dx \\
	&= \int_{-\infty}^{+\infty}\left\{I_{j, n_1-j+1} \circ I_{k, n_2-k+1}^{-1} \circ I_{ r_{2,k}, N_2- r_{2,k}+1} \circ F_X(x) - I_{ r_{1,j}, N_1- r_{1,j}+1} \circ F_X(x)\right\}^+ \  d \ I_{ r_{2,k}, N_2- r_{2,k}+1}(F_X(x)) \\
	&= \int_0^1 \left\{I_{j, n_1-j+1} \circ I_{k, n_2-k+1}^{-1}(t) - I_{ r_{1,j}, N_1- r_{1,j}+1} \circ I_{ r_{2,k}, N_2- r_{2,k}+1}^{-1}(t)\right\}^+	 \ dt.
\end{align*}
This representation holds without additional assumptions. The positive-part operator appears because the integrand may change sign. Under condition~\eqref{eq:technical-condition}, the integrand is nonnegative, hence
\begin{align*}
	&\int_0^1 \left\{I_{j, n_1-j+1} \circ I_{k, n_2-k+1}^{-1}(t) - I_{ r_{1,j}, N_1- r_{1,j}+1} \circ I_{ r_{2,k}, N_2- r_{2,k}+1}^{-1}(t)\right\}^+\ dt \\
	&= \int_0^1 \left\{I_{j, n_1-j+1} \circ I_{k, n_2-k+1}^{-1}(t) - I_{ r_{1,j}, N_1- r_{1,j}+1} \circ I_{ r_{2,k}, N_2- r_{2,k}+1}^{-1}(t)\right\}\ dt.
\end{align*}
Thus,
\begin{align*}	
&\int_0^1 \left\{I_{j, n_1-j+1} \circ I_{k, n_2-k+1}^{-1}(t) - I_{ r_{1,j}, N_1- r_{1,j}+1} \circ I_{ r_{2,k}, N_2- r_{2,k}+1}^{-1}(t)\right\}^+\ dt \\
	&\qquad= \int_0^1 I_{j, n_1 - j + 1}(x)\, dI_{k, n_2 - k + 1}(x)
- \int_0^1 I_{r_{1,j}, N_1 - r_{1,j} + 1}(x)\, dI_{r_{2,k}, N_2 - r_{2,k} + 1}(x)\\
&\qquad= \int_0^1 I_{j, n_1 - j + 1}(x)\,
\frac{x^{k-1}(1-x)^{n_2-k}}{B(k, n_2-k+1)} \, dx \\
&\qquad\qquad- \int_0^1 I_{r_{1,j}, N_1-r_{1,j}+1}(x)\,
\frac{x^{r_{2,k}-1}(1-x)^{N_2-r_{2,k}}}{B(r_{2,k}, N_2-r_{2,k}+1)} \, dx\\
&\qquad = \int_0^1 \sum_{p=j}^{n_1}\binom{n_1}{p}
x^p(1-x)^{n_1-p}
\frac{x^{k-1}(1-x)^{n_2-k}}{B(k, n_2-k+1)} \, dx \\
&\qquad\qquad - \int_0^1 \sum_{p=r_{1,j}}^{N_1}\binom{N_1}{p}
x^p(1-x)^{N_1-p}
\frac{x^{r_{2,k}-1}(1-x)^{N_2-r_{2,k}}}{B(r_{2,k}, N_2-r_{2,k}+1)} \, dx\\ 
&\qquad = \sum_{p=j}^{n_1}\binom{n_1}{p}
\frac{B(k+p, n_2+n_1-p-k+1)}{B(k, n_2-k+1)} \\
&\qquad\qquad - \sum_{p=r_{1,j}}^{N_1}\binom{N_1}{p}
\frac{B(r_{2,k}+p, N_2+N_1-p-r_{2,k}+1)}{B(r_{2,k}, N_2-r_{2,k}+1)}.
\end{align*}
\end{proof}



%%%% Proof of Proposition 3.5
\subsection{Proof of Proposition~\ref{prop:condition}}

\begin{proof}
Specification of the selection rule $S_{i,j} = 1 \;\Longleftrightarrow\; X_{i,j} \geq X_{i,(N_i - n_i + 1)}$ implies that $r_{i,j} = N_i + 1 - j$. When either $n_1 = N_1$ and $n_2 = N_2$ or $n_1 = n_2$ and $N_1 = N_2$ we have, for all $t \in (0,1)$, that
$$
  I_{N_2-j,j+1} \circ I_{N_1-j,j+1}^{-1} \circ I_{n_1-j,j+1} \circ I_{n_2-j,j+1}^{-1} (t) \geq t.
$$
The left hand side of the above equation increases whenever $n_1$ or $N_1$ are increased. This follows from the monotonicity of the inverse incomplete beta function in its parameters \citep{askitis2021logarithmic}, together with the fact that stochastic dominance is preserved under monotone transformations. Therefore,
$
  I_{N_2-j,j+1} \circ I_{N_1-j,j+1}^{-1} \circ I_{n_1-j,j+1} \circ I_{n_2-j,j+1}^{-1} (t) \geq t
$
whenever $\frac{n_1}{n_2}\leq 1$ and $\frac{N_2}{N_1} \leq 1$. Observe that
\begin{align*}
&I_{N_2-j,j+1} \circ I_{N_1-j,j+1}^{-1} \circ I_{n_1-j,j+1} \circ I_{n_2-j,j+1}^{-1} (t) \geq t, \\ 
&\quad\implies I_{n_1-j,j+1} \circ I_{n_2-j,j+1}^{-1} (t) \geq I_{N_1-j,j+1} \circ I_{N_2-j,j+1}^{-1}(t),   \\
&\quad\implies I_{n_1-j,j+1} \circ I_{n_2-j,j+1}^{-1} (x) - I_{N_1-j,j+1} \circ I_{N_2-j,j+1}^{-1}(t) \geq t.
\end{align*}
Technical condition \eqref{eq:technical-condition} holds and the conclusion of Theorem~\ref{theorem:samedistgeneral} applies.
\end{proof}



%%%% Proof of Proposition 4.2
\subsection{Proof of Proposition~\ref{prop:weakstodom}}

\begin{proof}
First, in combination with the Definition \ref{def:stodom} and the assumption $\frac{n_1}{N_1} \leq \frac{n_2}{N_2}$, we have 
$$
  F_{Y_1}(x)^{\frac{n_1}{N_1}} \geq F_{Y_2}(x)^{\frac{n_2}{N_2}} 
 \implies  
 x \geq F_{Y_1}^{-1}\left(F_{Y_2}(x)^{\frac{n_2N_1}{N_2n_1}}\right).
$$ 
By applying the Equation \ref{eq:pararep2} from the Lemma \ref{lemma:rep} we have 
\begin{align*}
	&\Prob\left(X_{1,(N_1)} > X_{2,(N_2)}, Y_{1,(n_1)} < Y_{2,(n_2)}\right) 
	= \Prob\left( g_{n_1,n_2}\left(Y_{2,(n_2)}\right) < Y_{1,(n_1)}  < Y_{2,(n_2)}\right) \\
	&\qquad = \Prob\left( F_{Y_1}^{-1}\left(F_{Y_2}(Y_{2,(n_2)})^{\frac{n_2N_1}{N_2n_1}}\right) < Y_{1,(n_1)}  < Y_{2,(n_2)}\right) \\
	&\qquad = \int_{-\infty}^{+\infty} \int_{F_{Y_1}^{-1}\left(F_{Y_2}(x)^{\frac{n_2N_1}{N_2n_1}}\right)}^{x} d F_{Y_1, (n_1)}(y) d F_{Y_2, (n_2)}(x) \\
	&\qquad = \int_{-\infty}^{+\infty} F_{Y_1}(x)^{n_1} - F_{Y_2}(x)^{\frac{n_2N_1}{N_2}}d F_{Y_2, (n_2)}(x)\\
	&\qquad = \int_0^1 F_{Y_1}\left( F_{Y_2}^{-1}(t^{\frac{1}{n_2}}) \right)^{n_1} - t^{\frac{N_1}{N_2}} dt. \\
\end{align*}
Based on the Definition \ref{def:stodom}, we have
\begin{align*}
	&\Prob\left(X_{1,(N_1)} > X_{2,(N_2)}, Y_{1,(n_1)} < Y_{2,(n_2)}\right) 
	= \int_0^1 F_{Y_1}\left( F_{Y_2}^{-1}(t^{\frac{1}{n_2}}) \right)^{n_1} - t^{\frac{N_1}{N_2}} dt \\
	&\qquad \geq \int_0^1 F_{Y_2}\left( F_{Y_2}^{-1}(t^{\frac{1}{n_2}}) \right)^{n_1} - t^{\frac{N_1}{N_2}} dt \\
	&\qquad = \frac{n_2}{n_1+n_2} - \frac{N_2}{N_1+N_2},
\end{align*}
and
\begin{align*}
	&\Prob\left(X_{1,(N_1)} > X_{2,(N_2)}, Y_{1,(n_1)} < Y_{2,(n_2)}\right) 
	= \int_0^1 F_{Y_1}\left( F_{Y_2}^{-1}(t^{\frac{1}{n_2}}) \right)^{n_1} - t^{\frac{N_1}{N_2}} dt \\
	&\qquad \leq \int_0^1 \left(t^{1/n_2} + \varepsilon\right)^{n_1} - t^{\frac{N_1}{N_2}} dt\\
	&\qquad = n_2\sum_{i=0}^{n_1} \frac{{n_1 \choose i}\varepsilon^i}{n_1+n_2-i}- \frac{N_2}{N_1 + N_2}.
\end{align*}
\end{proof}



%%%% Proof of Theorem 4.3
\subsection{Proof of Theorem~\ref{theorem:normalgeneral}}

\begin{proof}
Equation \ref{eq:pararep2} from Lemma \ref{lemma:rep} yields
\begin{align*}
	g_{n_1, n_2}(x) = F_{Y_1}^{-1}\left(F_{Y_2}(x)^{\frac{n_2N_1}{N_2n_1}}\right) = \Phi^{-1}\left(\Phi\left(\frac{x-\mu_2}{\sigma_2}\right)^\frac{n_2N_1}{N_2n_1}\right)\sigma_1 +\mu_1.
\end{align*}
Let $A$ and $R$ be as defined, respectively, in \eqref{eq:A} and \eqref{eq:remainder}. We have,
\begin{align*}
  &\Prob\left(X_{1,(N_1)} > X_{2,(N_2)}, Y_{1,(n_1)} < Y_{2,(n_2)}\right) = \Prob\left( g_{n_1,n_2}\left(Y_{2,(n_2)}\right) < Y_{1,(n_1)}  < Y_{2,(n_2)}\right)\\
	&\qquad= \Prob\left( F_{Y_1}^{-1}\left(F_{Y_2}(Y_{2,(n_2)})^{\frac{n_2N_1}{N_2n_1}}\right) < Y_{1,(n_1)}  < Y_{2,(n_2)}\right)\\
	&\qquad= \int_A \int_{F_{Y_1}^{-1}\left(F_{Y_2}(x)^{\frac{n_2N_1}{N_2n_1}}\right)}^{x} d F_{Y_1, (n_1)}(y) d F_{Y_2, (n_2)}(x)\\	
	&\qquad= \int_{A} F_{Y_1}(x)^{n_1} - F_{Y_2}(x)^{\frac{n_2N_1}{N_2}}d F_{Y_2, (n_2)}(x) \\
	&\qquad= \int_{-\infty}^{\infty} F_{Y_1}(x)^{n_1} - F_{Y_2}(x)^{\frac{n_2N_1}{N_2}}d F_{Y_2, (n_2)}(x) - \int_{A^c} F_{Y_1}(x)^{n_1} - F_{Y_2}(x)^{\frac{n_2N_1}{N_2}}d F_{Y_2, (n_2)}(x) \\
	&\qquad= \int_{-\infty}^{\infty} F_{Y_1}(x)^{n_1} - F_{Y_2}(x)^{\frac{n_2N_1}{N_2}}d F_{Y_2, (n_2)}(x) + R.
\end{align*}
We continue with the left-hand side of the last equality,	
\begin{align*}
&\int_{-\infty}^{\infty} \left[ F_{Y_1}(x)^{n_1} - F_{Y_2}(x)^{\frac{n_2N_1}{N_2}} \right] d F_{Y_2,(n_2)}(x) 
= \int_0^1 \left[ F_{Y_1}\!\left( F_{Y_2}^{-1}(t^{1/n_2}) \right)^{n_1} - t^{N_1/N_2} \right] dt\\[6pt]
&\qquad= \int_{0}^{1} \left( \Phi\!\left(\frac{\Phi^{-1}(t^{1/n_2})\sigma_2 + \mu_2 - \mu_1}{\sigma_1}\right) \right)^{n_1} \, dt \;-\; \frac{N_2}{N_1 + N_2} \\[6pt]
&\qquad= \int_{-\infty}^{+\infty} 
\Phi\!\left(\frac{\sigma_2 x + \mu_2 - \mu_1}{\sigma_1}\right)^{n_1} 
\left( n_2 \Phi(x)^{n_2-1} \phi(x) \right)\, dx \;-\; \frac{N_2}{N_1 + N_2} \\[8pt]
&\qquad= n_2 \int_{-\infty}^{+\infty}
\mathbb{P}\!\left(
\begin{array}{l}
Z_{1,1} - \frac{\sigma_2}{\sigma_1}X \le \frac{\mu_2 - \mu_1}{\sigma_1},\, \dots,\,
Z_{1,n_1} - \frac{\sigma_2}{\sigma_1}X \le \frac{\mu_2 - \mu_1}{\sigma_1},\\[3pt]
Z_{2,1}-X \le 0,\, \dots,\, Z_{2,n_2-1}-X \le 0
\end{array}
\;\middle|\; X=x \right)
\phi(x)\, dx \;-\; \frac{N_2}{N_1 + N_2} \\[10pt]
&\qquad= n_2 \,
\mathbb{P}\!\left(
\begin{array}{l}
Z_{1,1} - \frac{\sigma_2}{\sigma_1}X \le \frac{\mu_2 - \mu_1}{\sigma_1},\, \dots,\,
Z_{1,n_1} - \frac{\sigma_2}{\sigma_1}X \le \frac{\mu_2 - \mu_1}{\sigma_1},\\[3pt]
Z_{2,1}-X \le 0,\, \dots,\, Z_{2,n_2-1}-X \le 0
\end{array}
\right)
\;-\; \frac{N_2}{N_1 + N_2} \\[10pt]
&\qquad= n_2\, F_{\mathbf{W}}\!\left( \mathbf{w} = (
\frac{\mu_2 - \mu_1}{\sigma_1},\dots,\frac{\mu_2 - \mu_1}{\sigma_1},
0,\dots,0)
\right) \;-\; \frac{N_2}{N_1 + N_2}.
\end{align*}

The distribution function $F_{\mathbf{W}}$ corresponds to random variable $\mathbf{W} \sim \mathcal{N}_{(n_1 + n_2 - 1)}(\boldsymbol{0}, \boldsymbol{\Sigma})$ where $\boldsymbol{\Sigma}$ is as defined in \eqref{eq:Sigma} and \eqref{eq:Sigmastructure}. The first $n_1$ elements of $\mathbf{w}$ is $\frac{\mu_2 - \mu_1}{\sigma_1}$ and the last $n_2-1$ elements are 0.

The bound on the remainder \eqref{eq:remainder-bound} follows from the argument in Section~\ref{subsec:normal}.
\end{proof}














%%%% Proof of Theorem 4.4
\subsection{Proof of Theorem \ref{theorem:asymptotic}}
\label{sec:proofplugin}
\begin{proof}

\begin{align*}
	\mathbf F_{\mathbf W}(\mathbf w) &= \mathbb{P}\!\left(
\begin{array}{l}
Z_{1} + \frac{\sigma_2}{\sigma_1}Z_0 \le \frac{\mu_2 - \mu_1}{\sigma_1},\, \dots,\,
Z_{n_1} + \frac{\sigma_2}{\sigma_1}Z_0 \le \frac{\mu_2 - \mu_1}{\sigma_1},\\[3pt]
Z_{n_1+1}+ Z_0 \le 0,\, \dots,\, Z_{n_1+n_2-1}+Z_0 \le 0
\end{array}
\right) \\
&= \mathbb{P}\!\left(
\begin{array}{l}
Z_{1} \le \frac{\mu_2 - \mu_1}{\sigma_1} - \frac{\sigma_2}{\sigma_1}Z_0,\, \dots,\,
Z_{n_1} \le \frac{\mu_2 - \mu_1}{\sigma_1} - \frac{\sigma_2}{\sigma_1}Z_0,\\[3pt]
Z_{n_1+1}\le - Z_0 ,\, \dots,\, Z_{n_1+n_2-1} \le -Z_0
\end{array}
\right) \\
\end{align*}

Define $M_{n_1} = \max\{Z_1, Z_2, \dots, Z_{n_1}\}$, $M_{n_2-1} = \max\{Z_{n_1+1}, Z_{n_1+2}, \dots, Z_{n_1+n_2-1}\}$, then the CDF of $\mathbf W$ essentially computes the joint probability:

$$P(M_{n_1} \le \frac{\mu_2 - \mu_1}{\sigma_1} - \frac{\sigma_2}{\sigma_1} Z_0 \;\;\text{and}\;\; M_{n_2-1} \le -Z_0)$$

From extreme value theory, as $n_1 \to \infty$, $M_{n_1}$ sharply concentrates around $\sqrt{2 \ln n_1}$ \citep{embrechts2013modelling}. 

Substitute this into the first inequality:

$$\sqrt{2 \ln n_1} \le \frac{\mu_2 - \mu_1}{\sigma_1} - \frac{\sigma_2}{\sigma_1} Z_0 \implies Z_0 \le \frac{\frac{\mu_2 - \mu_1}{\sigma_1} - \sqrt{2 \ln n_1}}{\frac{\sigma_2}{\sigma_1}}$$

As $n_1$ grows, $\frac{\frac{\mu_2 - \mu_1}{\sigma_1} - \sqrt{2 \ln n_1}}{\frac{\sigma_2}{\sigma_1}}$ becomes a very large negative number. This forces the random variable $Z_0$ deep into its left tail. 

Similarly, as $n_2 \to \infty$, $M_{n_2-1}$ sharply concentrates around $\sqrt{2 \ln (n_2-1)}$ \citep{embrechts2013modelling}. 

Substitute this into the second inequality:

$$\sqrt{2 \ln (n_2-1)} \le -Z_0 \implies Z_0 \le -\sqrt{2 \ln (n_2-1)}$$

As $n_2$ grows, $-\sqrt{2 \ln (n_2-1)}$ becomes a very large negative number. This forces the random variable $Z_0$ deep into its left tail. 

Then the CDF of $\mathbf W$ essentially computes the probability: 

$$
\mathbf{P}(Z_0 \le \frac{\frac{\mu_2 - \mu_1}{\sigma_1} - \sqrt{2 \ln n_1}}{\frac{\sigma_2}{\sigma_1}}, Z_0 \le -\sqrt{2 \ln (n_2-1)}),
$$
as $n_1, n_2 \to \infty. $ 

This probability is dominated entirely by the tail of $Z_0$. We can use the Mill's ratio \citep{mills1926table} and apply it to standard Gaussian tail approximation, $\Phi(-x) \sim \frac{\phi(x)}{x}$ as $x \to \infty$. 

Since $$\lim_{n_1, n_2 \to \infty} \frac{\ln n_2}{\ln n_1} > \left( \frac{\sigma_1}{\sigma_2} \right)^2,$$ the tail behavior of $Z_0$ is dominated by $n_2$, then

 
$$
F_{\mathbf{W}}(\mathbf{w}) \approx \Phi(-\sqrt{2 \ln (n_2-1)}) \sim \frac{\phi(\sqrt{2 \ln (n_2-1)})}{\sqrt{2 \ln (n_2-1)}} = \frac{1}{\sqrt{4\pi \ln (n_2-1)}} \left( \frac{1}{n_2-1} + o(\frac{1}{n_2}) \right)
$$

Let $\mathbf{W^*} \sim \mathcal{N}_{(n_1 + n_2 - 1)}(\boldsymbol{0}, \boldsymbol{\Sigma^*})$ and covariance matrix 
\begin{equation}\label{eq:Sigma}
\boldsymbol{\Sigma^*} = \begin{bmatrix}
	\Sigma_{11}^* & \Sigma_{12}^*\\
	\Sigma_{21}^* & \Sigma_{22}^*\\
\end{bmatrix},
\end{equation}
defined with the following structure \\
\begin{align}\label{eq:Sigmastructure}
\begin{split}
\Sigma_{11}^* = \begin{bmatrix}
\frac{\sigma_4^2}{\sigma_3^2}+1 & \frac{\sigma_4^2}{\sigma_3^2} & \hdots & \frac{\sigma_2^2}{\sigma_3^2} \\
\frac{\sigma_4^2}{\sigma_3^2} & \ddots & \ddots  & \vdots \\
\vdots & \ddots & \ddots & \frac{\sigma_4^2}{\sigma_3^2} \\
\frac{\sigma_4^2}{\sigma_3^2} & \dots & \frac{\sigma_4^2}{\sigma_3^2} & \frac{\sigma_4^2}{\sigma_3^2}+1 \\
\end{bmatrix}_{n_1 \times n_1} \;\;
&\Sigma_{12}^{T*} = \Sigma_{21}^* = \begin{bmatrix}
\frac{\sigma_4}{\sigma_3} & \hdots & \hdots & \frac{\sigma_4}{\sigma_3}\\
\vdots & \ddots &  & \vdots\\
\vdots &  & \ddots & \vdots\\
\frac{\sigma_4}{\sigma_3} & \hdots &\hdots & \frac{\sigma_4}{\sigma_3}\\
\end{bmatrix}_{(n_2-1) \times n_1} \\
&\Sigma_{22}^* = \begin{bmatrix}
2 & 1 & \hdots & 1\\
1 & \ddots & \ddots & \vdots\\
\vdots & \ddots & \ddots & 1\\
1 & \hdots & 1 & 2\\
\end{bmatrix}_{(n_2-1) \times (n_2-1)}.
\end{split}
\end{align}

and $\mathbf{w}^* = (\frac{\mu_4 - \mu_3}{\sigma_3}, \dots, \frac{\mu_4 - \mu_3}{\sigma_3}, 0, \dots, 0)$, where the first $n_1$ elements of $\mathbf{w}^*$ are $\frac{\mu_4 - \mu_3}{\sigma_3}$ and the last $n_2-1$ elements are 0.

Then the $\mathbf F_{\mathbf W^*}(\mathbf w^*)$ is

\begin{align*}
	\mathbf F_{\mathbf W^*}(\mathbf w^*) &= \mathbb{P}\!\left(
\begin{array}{l}
Z_{1} + \frac{\sigma_4}{\sigma_3}Z_0 \le \frac{\mu_4 - \mu_3}{\sigma_3},\, \dots,\,
Z_{n_1} + \frac{\sigma_4}{\sigma_3}Z_0 \le \frac{\mu_4 - \mu_3}{\sigma_3},\\[3pt]
Z_{n_1+1}+ Z_0 \le 0,\, \dots,\, Z_{n_1+n_2-1}+Z_0 \le 0
\end{array}
\right) \\
&= \mathbb{P}\!\left(
\begin{array}{l}
Z_{1} \le \frac{\mu_4 - \mu_1}{\sigma_3} - \frac{\sigma_4}{\sigma_3}Z_0,\, \dots,\,
Z_{n_1} \le \frac{\mu_4 - \mu_1}{\sigma_3} - \frac{\sigma_4}{\sigma_3}Z_0,\\[3pt]
Z_{n_1+1}\le - Z_0 ,\, \dots,\, Z_{n_1+n_2-1} \le -Z_0
\end{array}
\right) \\
\end{align*}

Since the $\mu_3$, $\mu_4$, $\sigma_3$, and $\sigma_4$ satisfy the same assumptions as $\mu_1$, $\mu_2$, $\sigma_1$, and $\sigma_2$ in the theorem statement, and applying the same technique for $\mathbf F_{\mathbf W}(\mathbf w)$, we have

$$
F_{\mathbf{W}^*}(\mathbf{w}^*) \approx \Phi(-\sqrt{2 \ln (n_2-1)}) \sim \frac{\phi(\sqrt{2 \ln (n_2-1)})}{\sqrt{2 \ln (n_2-1)}} = \frac{1}{\sqrt{4\pi \ln (n_2-1)}} \left( \frac{1}{n_2-1} + o(\frac{1}{n_2}) \right)
$$

Consider the difference between paradox probabilities using two sets of different parameters, and 

\begin{align*}
		&\Prob\left(X_{1,(N_1)} > X_{2,(N_2)}, Y_{1,(n_1)} < Y_{2,(n_2)} \mid \mu_3, \mu_4, \sigma_3, \sigma_4 \right) \\
		&\qquad- \Prob\left(X_{1,(N_1)} > X_{2,(N_2)}, Y_{1,(n_1)} < Y_{2,(n_2)} \mid \mu_1, \mu_2, \sigma_1, \sigma_2\right) \\ 
		&= n_2(F_{\mathbf{W^*}}(\mathbf{w^*}) - F_{\mathbf{W}}(\mathbf{w})) + (R^* - R)\to 0
\end{align*} as $n_1, n_2 \to \infty. $ 
where $R^* := R(\mu_3,\mu_4,\sigma_3,\sigma_4, n_1,n_2,N_1,N_2)$ and $R := R(\mu_1,\mu_2,\sigma_1,\sigma_2, n_1,n_2,N_1,N_2)$ are defined in (\ref{eq:remainder})
\end{proof}



\section{Simulations}

\subsection{Simulation Results for Theorem \ref{theorem:normalgeneral}}

We now show the simulation results of our Normal Distribution example under different configurations of $(\mu_1,\mu_2,\sigma_1,\sigma_2, n_1,n_2,N_1,N_2)$. We verify Theorem~\ref{theorem:normalgeneral} in the numerical simulation. Directly paradox probabilities from these theorems are compared with Monte Carlo paradox probabilities in a variety of settings. Each of these settings will involve comparisons of the observed statistics and latent aptitudes of the highest achievers, the pair $(Y_{i,(n_i)}, X_{i,(N_i)})$ across two populations.

The simulations will proceed as follows: Distributions for $Y_1$ and $Y_2$ are specified along with sample sizes $n_i$, $N_i$, $i = 1,2$. We will specify a Monte Carlo sample size of 2000. For iteration $m$ of the Monte Carlo sample procedure $Y_{1,j}$, $j=1,\ldots,n_1$ and $Y_{2,k}$, $k=1,\ldots,n_2$ are generated. Then $X_{i,(N_i)}$, $i =1,2$ is calculated using \eqref{eq:truepar} where $X\sim\text{Pareto}(\alpha=1.16)$. Then 
$
  1(X_{(N_1)} > X_{(N_2)}, Y_{(n_1)} < Y_{(n_2)})
$
is recorded. 

The theoretical results are computed from the Theorem~\ref{theorem:normalgeneral} using mvtnorm package \citep{R-mvtnorm} in R.

The results are tabulated in Table~\ref{Tab:simulationnormal}. These results show strong agreement between our theory and numerical simulation. They also reveal that paradox probabilities are non-trivial across these configurations. 

\begin{table}[H]
\centering
\begin{tabular}{rrrrrrrr|r|r}
\hline
$n_1$ & $n_2$ & $N_1$ & $N_2$ & $\mu_1$ & $\mu_2$ & $\sigma_1$ & $\sigma_2$ & Simulation Prob & Theoretical Prob \\
\hline
10 & 40 & 400 & 600 & -1 & 1 & 1.0 & 3 & 0.4000 & 0.4000099 \\
10 & 20 & 400 & 600 & 0 & 0 & 1.0 & 3 & 0.3970 & 0.3970305 \\
10 & 40 & 400 & 300 & -1 & 1 & 1.0 & 3 & 0.5715 & 0.5714174 \\
10 & 20 & 400 & 600 & 0 & 1 & 1.5 & 2 & 0.3645 & 0.3647783 \\
10 & 40 & 400 & 600 & 0 & 1 & 1.0 & 3 & 0.3995 & 0.3999245 \\
20 & 40 & 400 & 600 & 0 & 1 & 1.0 & 3 & 0.3995 & 0.3999803 \\
10 & 40 & 400 & 600 & -1 & 0 & 1.5 & 3 & 0.4000 & 0.3995149 \\
20 & 40 & 400 & 600 & 0 & 0 & 1.5 & 2 & 0.2970 & 0.2962431 \\
20 & 40 & 400 & 600 & 0 & 0 & 1.0 & 3 & 0.3750 & 0.3740938 \\
10 & 20 & 400 & 600 & -1 & 1 & 1.0 & 3 & 0.3990 & 0.4000059 \\
\hline
\end{tabular}
\caption{Simulation Results for Normal Example. }
\label{Tab:simulationnormal}
\end{table}






\bibliography{FHM.bib}





\begin{theorem}[Asymptotic stability under plug-in estimation]
\label{theorem:asymptotic-clean}
Assume a Full House Model as in Definition~\ref{def:fhm} and suppose that $Y_i \sim N(\mu_i,\sigma_i^2)$, $i=1,2$. Assume that the selection indicators satisfy $S_{i,(N_i)} = 1$ for $i=1,2$. Let the paradox probability be given by Theorem~\ref{theorem:normalgeneral}.

Let $\hat{\mu}_i$ and $\hat{\sigma}_i$ be estimators satisfying
\[
\hat{\mu}_i - \mu_i = O_P(n_i^{-1/2}), 
\qquad
\hat{\sigma}_i - \sigma_i = O_P(n_i^{-1/2}),
\quad i=1,2.
\]

Suppose that $n_1,n_2 \to \infty$ with
\[
\frac{\ln n_2}{\ln n_1} > \left(\frac{\sigma_1}{\sigma_2}\right)^2,
\]
and that $\frac{n_1}{N_1} \le \frac{n_2}{N_2}$, $\sigma_1 \le \sigma_2$, and $\mu_1 \le \mu_2$.

Then the plug-in estimator of the paradox probability is consistent, i.e.
\[
P(\hat{\theta}) - P(\theta) \xrightarrow{P} 0,
\]
where $\theta = (\mu_1,\mu_2,\sigma_1,\sigma_2)$.
\end{theorem}

\subsection{Proof of Theorem~\ref{theorem:asymptotic-clean}}

\begin{proof}
By Theorem~\ref{theorem:normalgeneral}, the paradox probability admits the decomposition
\[
P(\theta)
=
n_2 F_{\mathbf W}(\mathbf w(\theta)) - \frac{N_2}{N_1+N_2} + R(\theta),
\]
where $R(\theta)$ is the remainder term defined in \eqref{eq:remainder}. Let $\hat{\theta}$ denote the plug-in estimators. Then
\[
P(\hat{\theta}) - P(\theta)
=
n_2\Big(F_{\mathbf W}(\mathbf w(\hat{\theta})) - F_{\mathbf W}(\mathbf w(\theta))\Big)
+ \big(R(\hat{\theta}) - R(\theta)\big).
\]
We analyze the two terms separately.

\medskip

\noindent\textbf{Step 1: Control of the remainder term.}

From \eqref{eq:remainder-bound}, under the stated assumptions,
\[
|R(\theta)| \le \left(\frac{1}{2}\right)^{n_2}.
\]
Since this bound depends only on $(n_1,n_2)$, it also holds for $R(\hat{\theta})$ for sufficiently large samples. Therefore,
\[
|R(\hat{\theta}) - R(\theta)|
\le 2\left(\frac{1}{2}\right)^{n_2}
\;\longrightarrow\; 0.
\]

\medskip

\noindent\textbf{Step 2: Reduction to a Gaussian tail event.}

From the representation of $F_{\mathbf W}$ established in the proof of Theorem~\ref{theorem:normalgeneral}, we may write
\[
F_{\mathbf W}(\mathbf w(\theta))
=
\Prob\!\left(
M_{n_1} \le \frac{\mu_2 - \mu_1}{\sigma_1} - \frac{\sigma_2}{\sigma_1} Z_0,\;
M_{n_2-1} \le -Z_0
\right),
\]
where $Z_0, Z_1, \ldots$ are independent standard normal random variables, and
\[
M_{n_1} = \max_{1 \le j \le n_1} Z_j, 
\qquad
M_{n_2-1} = \max_{1 \le j \le n_2-1} Z_{n_1+j}.
\]

Let $a_n = \sqrt{2\ln n}$. It is well known that for any fixed $C>0$,
\[
\Prob\big(|M_n - a_n| > C\big) \to 0.
\]
Hence,
\[
M_{n_1} = a_{n_1} + O_P(1), 
\qquad
M_{n_2-1} = a_{n_2} + O_P(1).
\]

Define
\begin{align*}
A_n 
&:= \left\{
M_{n_1} \le \frac{\mu_2 - \mu_1}{\sigma_1} - \frac{\sigma_2}{\sigma_1} Z_0,\;
M_{n_2-1} \le -Z_0
\right\},\\
B_n 
&:= \left\{
a_{n_1} \le \frac{\mu_2 - \mu_1}{\sigma_1} - \frac{\sigma_2}{\sigma_1} Z_0 + C,\;
a_{n_2} \le -Z_0 + C
\right\}.
\end{align*}

By concentration of maxima,
\[
\Prob(A_n \triangle B_n) \to 0,
\]
and therefore $\Prob(A_n) = \Prob(B_n) + o(1)$.

The two constraints in $B_n$ can be written as
\[
Z_0 \le \frac{\frac{\mu_2 - \mu_1}{\sigma_1} - a_{n_1} + C}{\sigma_2/\sigma_1},
\qquad
Z_0 \le -a_{n_2} + C.
\]

Under the growth condition
\[
\frac{\ln n_2}{\ln n_1} > \left(\frac{\sigma_1}{\sigma_2}\right)^2,
\]
we have
\[
a_{n_1}\frac{\sigma_1}{\sigma_2} = o(a_{n_2}),
\]
so the first constraint is asymptotically weaker. Hence
\[
\Prob(B_n \triangle \{Z_0 \le -a_{n_2} + C\}) \to 0.
\]

It follows that
\[
F_{\mathbf W}(\mathbf w(\theta))
=
\Prob(Z_0 \le -\sqrt{2\ln n_2} + C) + o(1).
\]

Using standard Gaussian tail asymptotics, uniformly for fixed $C$,
\[
\Prob(Z_0 \le -x + C) = \Phi(-x)(1 + o(1)),
\]
we conclude that
\[
F_{\mathbf W}(\mathbf w(\theta))
=
\Phi(-\sqrt{2\ln n_2})(1 + o(1)).
\]

\medskip

\noindent\textbf{Step 3: Stability under parameter perturbations.}

The mapping $\theta \mapsto \mathbf w(\theta)$ is smooth, so
\[
\|\mathbf w(\hat{\theta}) - \mathbf w(\theta)\|
=
O_P(n_1^{-1/2} + n_2^{-1/2}).
\]

From Step 2, $F_{\mathbf W}(\mathbf w(\theta))$ is asymptotically equivalent to a one-dimensional Gaussian tail probability with threshold $x = \sqrt{2\ln n_2}$. Small perturbations of $\theta$ induce perturbations $\delta = O_P(n_1^{-1/2} + n_2^{-1/2})$ in this threshold.

Since $n_i^{-1/2} = o(\sqrt{\ln n_2})$, we may apply the expansion
\[
\Phi(-x + \delta) - \Phi(-x)
=
\phi(x)\delta + o(\delta \phi(x)),
\quad x \to \infty.
\]

Thus,
\[
F_{\mathbf W}(\mathbf w(\hat{\theta})) - F_{\mathbf W}(\mathbf w(\theta))
=
O_P\!\left(
\phi(\sqrt{2\ln n_2}) (n_1^{-1/2} + n_2^{-1/2})
\right).
\]

\medskip

\noindent\textbf{Step 4: Tail evaluation.}

Using
\[
\phi(\sqrt{2\ln n_2})
=
\frac{1}{\sqrt{2\pi}} e^{-\ln n_2}
=
O\!\left(\frac{1}{n_2}\right),
\]
we obtain
\[
F_{\mathbf W}(\mathbf w(\hat{\theta})) - F_{\mathbf W}(\mathbf w(\theta))
=
O_P\!\left(
\frac{1}{n_2}(n_1^{-1/2} + n_2^{-1/2})
\right).
\]

Multiplying by $n_2$,
\[
n_2\Big(F_{\mathbf W}(\mathbf w(\hat{\theta})) - F_{\mathbf W}(\mathbf w(\theta))\Big)
=
O_P(n_1^{-1/2} + n_2^{-1/2})
\;\longrightarrow\; 0.
\]

\medskip

\noindent\textbf{Step 5: Conclusion.}

Combining Steps 1 and 4 yields
\[
P(\hat{\theta}) - P(\theta) \xrightarrow{P} 0,
\]
as claimed.
\end{proof}



\end{document}
